\section{Symbolic and numerical validation}
\label{sec:validation}

The accompanying script \texttt{validation/verify\_geometry.py} performs two
levels of checks. Symbolically it verifies:
\begin{itemize}
 \item \(\matr Q_U=\matr A^{\mathsf T}\matr A\), the voltage null vector,
       and the torque-matrix eigenvalues;
 \item the PSM center, its conic trajectory, and the double-angle formula;
 \item both inverse coordinate transformations, canonical torque, and loss;
 \item the hyperbolic parameterization, projective \(J\)-circles, and complex
       square map;
 \item the explicit EESM capability coefficient \(\kappa\), its circular
       machine-design boundary, and the \(R_{\mathrm{exc}}\downarrow0\) limit;
 \item nonsalient and salient PSM loss-limited maxima, the characteristic-
       current design window, and multi-point set intersection logic.
\end{itemize}
Random regression draws positive resistances and inductances, signed speeds,
and currents. For every draw it evaluates the original and transformed
expressions for loss, torque, and squared voltage and requires relative errors
below \(10^{-10}\). It also checks feasibility reconstruction from performance
coordinates.

The normalized toy case
\begin{equation}
 R_s=0.4,\quad R_{\mathrm{exc}}=0.7,\quad L_d=1.2,\quad L_q=0.8,
 \quad L_{\mathrm{exc}}=0.9,\quad p=2
 \label{eq:toy-case}
\end{equation}
is dimensionless and selected to avoid accidental symmetry. It is not machine
data. The engineering-scale PSM values used only to exercise plotting are also
marked synthetic. The workspace contains no traceable measurement or project
dataset for this machine problem; unrelated CSV and MATLAB files belong to
other publications and were not repurposed.

\begin{table}[tb]
\centering
\caption{Independent validation routes and their failure signatures.}
\label{tab:validation}
\begin{tabularx}{\linewidth}{@{}p{.27\linewidth}X X@{}}
\toprule
Claim & Independent check & Typical detected error\\
\midrule
Voltage cylinder & rank/SVD and explicit null vector & sign or column error in \(\matr A\)\\
PSM center & square completion and zero-voltage solve & incorrect PM affine term\\
Whitening & quadratic congruence and random values & wrong resistance square root\\
Torque rotation & orthogonality and inverse map & swapped active/neutral signs\\
Projective circles & symbolic completion and sampled levels & missing \(1/\eta\) factor\\
Performance frontier & sampled sphere and support eigenvectors & wrong eigenvalue branch\\
Inverse EESM map & direct generalized eigenvalue and parameter sweeps & missing loss factor or false dimensionless claim\\
PSM design window & angular maximization and inequality endpoints & sufficient field-weakening test reported as universal\\
Optimum & primal, complementarity, and KKT residuals & infeasible stationary candidate\\
\bottomrule
\end{tabularx}
\end{table}

Plots are evidence only after residuals pass. A finite plotting box can make
the voltage cylinder appear closed and can create a false maximum-torque point
on a hidden boundary. Every numerical optimum therefore reports active
physical constraints and is rerun with enlarged diagnostic bounds.

The dynamic extension has an independent script,
\texttt{validation/verify\_dynamics.py}. It verifies symmetry and positive
definiteness of the reciprocal inductance matrix, the dynamic-whitening
identity \(\matr D\matr B\matr S^{1/2}=\matr I\),
\(\matr G=\matr D^{\mathsf T}\matr D\), the metric tangent projector, and
positive finite-horizon Gramian eigenvalues. Every reported piecewise-linear
path is resampled for stator-circle and field-voltage utilization. Because the
required voltage is affine along one segment, convexity proves that its
maximum norm and maximum field-voltage magnitude occur at a segment endpoint;
the dense resampling remains as a regression check. The six trajectory files
under \texttt{figures/data} are regenerated from the same declared model used
in Table~\ref{tab:dynamic-results}.
